\section*{Chapitre \ref{ch:g5:raw-materials-and-recycling} --- Des matières premières au recyclage}

\begin{solution}{exo:g5:raw-materials-and-recycling:1}
Naturels~: le bois, la laine, la pierre. Fabriqués~: le verre, l'acier,
le plastique.
\end{solution}

\begin{solution}{exo:g5:raw-materials-and-recycling:2}
Le papier~: le bois (les arbres). Le verre~: le sable (avec de la soude
et du calcaire). Le coton~: le cotonnier. Le plastique~: le pétrole
brut.
\end{solution}

\begin{solution}{exo:g5:raw-materials-and-recycling:3}
Une roche dont on extrait un métal. L'aluminium vient de la bauxite.
\end{solution}

\begin{solution}{exo:g5:raw-materials-and-recycling:4}
Le bocal dans le bac à verre, le journal avec le papier et le carton, la
canette avec le métal, la bouteille avec les bouteilles en plastique.
\end{solution}

\begin{solution}{exo:g5:raw-materials-and-recycling:5}
La \omterm{def:g5:raw-materials-and-recycling:raw-material}{matière première} neuve~: le matériau recyclé est utilisé à la place
du sable, du \omterm{def:g5:raw-materials-and-recycling:raw-material}{minerai} ou du pétrole prélevés dans la nature.
\end{solution}

\begin{solution}{exo:g5:raw-materials-and-recycling:6}
Renouvelables~: le bois (si l'on replante), la laine, le coton. Non
renouvelables~: le pétrole brut, le \omterm{def:g5:raw-materials-and-recycling:raw-material}{minerai} de fer.
\end{solution}

\begin{solution}{exo:g5:raw-materials-and-recycling:7}
Fabriquer du verre à partir de sable est une \omterm{def:g5:irreversible-changes:chemical-change}{transformation chimique}~:
quelque chose de nouveau est fabriqué, et il n'y a pas de chemin de
retour simple.
\end{solution}

\begin{solution}{exo:g5:raw-materials-and-recycling:8}
Par exemple~: réduire --- acheter des fruits sans emballage en
plastique~; réutiliser --- remplir de nouveau un bocal en verre~;
recycler --- mettre une canette dans le bac à métal.
\end{solution}

\begin{solution}{exo:g5:raw-materials-and-recycling:9}
Chaque matériau se recycle à sa façon~: des morceaux d'un autre
matériau gâcheraient le verre fondu.
\end{solution}

\begin{solution}{exo:g5:raw-materials-and-recycling:10}
Un vingtième de 20 unités~: $20 \div 20 = 1$ unité d'énergie, environ.
\end{solution}

\begin{solution}{exo:g5:raw-materials-and-recycling:11}
Le \omterm{def:g5:raw-materials-and-recycling:recycling}{recyclage} demande encore de l'énergie, des machines et des
transports~; un objet qui n'est jamais fabriqué économise sa \omterm{def:g5:raw-materials-and-recycling:raw-material}{matière
première} et toute cette énergie.
\end{solution}
